\section{Introduction}
\label{sec:intro}

LLM agents increasingly acquire specialised abilities from \emph{skills}: directories holding
instructions, scripts and reference files that the agent loads when a task calls for them
\citep{anthropic2025skillsblog,agentskills2026spec}. A good skill encodes expertise that took time
and money to produce, such as how a tax authority cross-checks a filing or how a firm scores a
contract, and skills are starting to be sold. They ship as plain text. Whoever obtains the folder
can read it, republish it, or load it into a competing product, and the licence terms attached to
it enforce nothing.

This paper asks what it would take to sell a skill that only a chosen AI provider can open, and
how much protection that actually buys. We make three observations. First, the cryptography is
straightforward and cheap: an owner can encrypt a skill to a public key the provider publishes, in
the same way anyone can encrypt email to a published key, without the parties ever talking.
Second, the protection depends on a strict architectural rule that is easy to get wrong: the skill
must only ever be decrypted on the provider's servers, never in a client, because anything present
on the user's device can be extracted by the user. Third, and most important, encryption cannot
close the channel that matters most. To follow a skill, the model must read it in clear text, and
the model's answers go to the user, who can ask for the skill directly, indirectly, or in
disguise.

\paragraph{Contributions.}
\begin{itemize}
  \item \textbf{Design.} \sys{}, a provider-bound packaging scheme for agent skills: HPKE key
  wrapping to a key published at \texttt{/.well-known/sealed-skill-keys}, STREAM-encrypted
  payloads, owner signatures, offline-verifiable licences, and per-customer canaries, with no
  owner--provider handshake (\cref{sec:design}). We state the server-side decryption rule and why
  it is necessary (\cref{sec:boundary}), and analyse security against five adversaries
  (\cref{sec:security}).
  \item \textbf{Implementation.} An open prototype with a mock provider, and measurements showing
  sealing costs are negligible next to model inference (\cref{sec:rq1}).
  \item \textbf{Measurement of the residual channel.} On eight synthetic multi-rule skills, 24
  attacks and two open models, we quantify extraction from context, the effect and utility cost of
  a guard instruction, and the trade-off curves of lexical and semantic output filters
  (\cref{sec:rq2,sec:rq3}). We use an evidence-grounded judge for rule disclosure after finding that
  an unconstrained LLM judge counts applying a rule as disclosing it.
  \item \textbf{Scripts.} A tool-use case study showing that ``run-only'' scripts cannot be
  enforced in a sandbox that also executes arbitrary code, and that owner-hosted tools protect
  implementations but not function (\cref{sec:casestudy}).
\end{itemize}

\paragraph{Findings in brief.}
Sealing removes every non-model path to a skill at negligible cost. Once a skill is in context,
though, most of it can be extracted: without defences, 72--75\% of attack conversations disclosed
at least half of a skill's rules. A guard instruction and a de-obfuscating lexical filter together
reduce this to about 2\%, but block 12.5--19\% of legitimate replies, because models quote rules
while using them. Sealing is therefore protection against copying and redistribution plus a
measurable reduction in extraction, not secrecy. For content that must stay secret, the logic has
to leave the model's context entirely, and even then its function remains observable, as recent
skill-stealing work demonstrates~\citep{tsai2026daydreaming,hua2026behavioral}.

The closest prior design, by \citet{vanwyk2023protect}, decrypts an encrypted prompt through a
bridge placed ideally at the provider; we differ in using a published provider key with no
handshake or bridge operator, in protecting whole skills including scripts, and in measuring the
extraction channel (\cref{sec:related}).
